% GENERATED FILE. Edit canonical lessons, then run npm run generate:books.
% canonical-source-hash: fnv1a64:526c6c20406b1e25
% canonical-lessons: JA-C287-yochiyochi, JA-C287-haeru, JA-C287-obaasan, JA-C287-shinseki, JA-C287-seichou-suru

\chapter{Growing Up}
\label{ch:ja-growing-up}
\hlchaptermodality{287}

\begin{chapteropening}
\textbf{By the end of this chapter:} \emph{I can say how a child grew and who came to see.}

Complete the last lesson of chapter 287: I can say how a child grew and who came to see.
\end{chapteropening}

\section[\texorpdfstring{\ja{よちよち} (yochiyochi)}{よちよち (yochiyochi)}]{\ja{よちよち} (yochiyochi) — toddling; with tottering steps}

\label{lesson:JA-C287-yochiyochi}



\begin{quote}
Before the new one: say the Japanese for congratulations, then the Japanese for a grandchild.
\end{quote}

\subsection*{You'll want to know: \ja{よちよち}}
\textbf{\ja{よちよち}} — \emph{yochiyochi} — \textquotedblleft{}toddling; with tottering steps\textquotedblright{}.

\begin{quote}
\textbf{\ja{あかちゃんが} \ja{よちよち} \ja{あるいています。}} — \emph{akachan ga yochiyochi aruite imasu} — \textquotedblleft{}The baby is toddling along.\textquotedblright{}
\end{quote}

\begin{cousinweb}
\textbf{\ja{よちよち}} (\emph{yochiyochi}) is a native Japanese mimetic word, inherited rather than borrowed from Chinese, and its older origin is not settled (JMdict).
\end{cousinweb}

\begin{grammarlens}[title={What you've built}]
One more word for telling what happened.
\end{grammarlens}

\subsection*{Your turn}
\begin{itemize}
\raggedright
  \item \emph{Say it:} \emph{yochiyochi}
  \item \emph{Say it:} \emph{yochiyochi}, once more
  \item \emph{Say it:} \emph{akachan ga yochiyochi aruite imasu} — The baby is toddling along.
  \item \emph{Recall:} say the Japanese for congratulations, then the Japanese for a grandchild, then say \emph{yochiyochi} again
\end{itemize}

\begin{tcolorbox}[breakable,colback=teal!4,colframe=teal!35!black,title={Before you move on}]
What does \ja{よちよち} mean? (\textquotedblleft{}Toddling; with tottering steps\textquotedblright{}.) Say it once more.
\end{tcolorbox}
\section[\texorpdfstring{\ja{はえる} (haeru)}{はえる (haeru)}]{\ja{はえる} (haeru) — to sprout; to come through (of teeth)}

\label{lesson:JA-C287-haeru}



\begin{quote}
Before the new one: say the Japanese for a grandchild, then the Japanese for toddling.
\end{quote}

\subsection*{You'll want to know: \ja{はえる}}
\textbf{\ja{はえる}} — \emph{haeru} — \textquotedblleft{}to sprout; to come through (of teeth)\textquotedblright{}.

\begin{quote}
\textbf{\ja{にわに} \ja{くさが} \ja{はえています。}} — \emph{niwa ni kusa ga haete imasu} — \textquotedblleft{}Grass is coming up in the garden.\textquotedblright{}
\end{quote}

Polite \textbf{\ja{はえます}}, \emph{haemasu}; an ichidan verb. Past (polite): \textbf{\ja{あかちゃんの} \ja{はが} \ja{はえました。}}, \emph{akachan no ha ga haemashita}.

\begin{cousinweb}
\textbf{\ja{はえる}} (\emph{haeru}) is a native Japanese word (wago), inherited rather than borrowed from Chinese; its kanji was attached later as a kun'yomi (JMdict).
\end{cousinweb}

\begin{grammarlens}[title={What you've built}]
One more word for telling what happened.
\end{grammarlens}

\subsection*{Your turn}
\begin{itemize}
\raggedright
  \item \emph{Say it:} \emph{haeru}
  \item \emph{Say it:} \emph{haeru}, once more
  \item \emph{Say it:} \emph{niwa ni kusa ga haete imasu} — Grass is coming up in the garden.
  \item \emph{Recall:} say the Japanese for a grandchild, then the Japanese for toddling, then say \emph{haeru} again
\end{itemize}

\begin{tcolorbox}[breakable,colback=teal!4,colframe=teal!35!black,title={Before you move on}]
What does \ja{はえる} mean? (\textquotedblleft{}To sprout; to come through (of teeth)\textquotedblright{}.) Say it once more.
\end{tcolorbox}
\section[\texorpdfstring{\ja{おばあさん} (obāsan)}{おばあさん (obāsan)}]{\ja{おばあさん} (obāsan) — a grandmother; an old woman}

\label{lesson:JA-C287-obaasan}



\begin{quote}
Before the new one: say the Japanese for toddling, then the Japanese for to sprout.
\end{quote}

\subsection*{You'll want to know: \ja{おばあさん}}
\textbf{\ja{おばあさん}} — \emph{obāsan} — \textquotedblleft{}a grandmother; an old woman\textquotedblright{}.

\begin{quote}
\textbf{\ja{おばあさんは} \ja{まいにち} \ja{さんぽします。}} — \emph{obāsan wa mainichi sanpo shimasu} — \textquotedblleft{}Grandma goes for a walk every day.\textquotedblright{}
\end{quote}

\begin{cousinweb}
\textbf{\ja{おばあさん}} (\emph{obāsan}) is a native Japanese word (wago), inherited rather than borrowed from Chinese: \emph{bāsan}, \textquotedblleft{}an old woman\textquotedblright{}, with the polite prefix \emph{o-}. JMdict writes it with characters for \textquotedblleft{}grandmother\textquotedblright{}, read for their meaning (JMdict).
\end{cousinweb}

\begin{grammarlens}[title={What you've built}]
One more word for telling what happened.
\end{grammarlens}

\subsection*{Your turn}
\begin{itemize}
\raggedright
  \item \emph{Say it:} \emph{obāsan}
  \item \emph{Say it:} \emph{obāsan}, once more
  \item \emph{Say it:} \emph{obāsan wa mainichi sanpo shimasu} — Grandma goes for a walk every day.
  \item \emph{Recall:} say the Japanese for toddling, then the Japanese for to sprout, then say \emph{obāsan} again
\end{itemize}

\begin{tcolorbox}[breakable,colback=teal!4,colframe=teal!35!black,title={Before you move on}]
What does \ja{おばあさん} mean? (\textquotedblleft{}A grandmother; an old woman\textquotedblright{}.) Say it once more.
\end{tcolorbox}
\section[\texorpdfstring{\ja{しんせき} (shinseki)}{しんせき (shinseki)}]{\ja{しんせき} (shinseki) — a relative}

\label{lesson:JA-C287-shinseki}



\begin{quote}
Before the new one: say the Japanese for to sprout, then the Japanese for a grandmother.
\end{quote}

\subsection*{You'll want to know: \ja{しんせき}}
\textbf{\ja{しんせき}} — \emph{shinseki} — \textquotedblleft{}a relative\textquotedblright{}.

\begin{quote}
\textbf{\ja{しょうがつに} \ja{しんせきが} \ja{あつまります。}} — \emph{shōgatsu ni shinseki ga atsumarimasu} — \textquotedblleft{}Our relatives gather at the new year.\textquotedblright{}
\end{quote}

\begin{cousinweb}
\textbf{\ja{しんせき}} (\emph{shinseki}) is Sino-Japanese (kango): two Chinese characters read with their on'yomi, readings borrowed from Chinese. They are \emph{shin}, \textquotedblleft{}relative\textquotedblright{} (Middle Chinese \emph{tshin}) and \emph{seki}, \textquotedblleft{}relatives\textquotedblright{} (\emph{tshek}) (KANJIDIC2; Middle Chinese in Baxter's transcription).
\end{cousinweb}

\begin{grammarlens}[title={What you've built}]
One more word for telling what happened.
\end{grammarlens}

\subsection*{Your turn}
\begin{itemize}
\raggedright
  \item \emph{Say it:} \emph{shinseki}
  \item \emph{Say it:} \emph{shinseki}, once more
  \item \emph{Say it:} \emph{shōgatsu ni shinseki ga atsumarimasu} — Our relatives gather at the new year.
  \item \emph{Recall:} say the Japanese for to sprout, then the Japanese for a grandmother, then say \emph{shinseki} again
\end{itemize}

\begin{tcolorbox}[breakable,colback=teal!4,colframe=teal!35!black,title={Before you move on}]
What does \ja{しんせき} mean? (\textquotedblleft{}A relative\textquotedblright{}.) Say it once more.
\end{tcolorbox}
\section[\texorpdfstring{\ja{せいちょうする} (seichō suru)}{せいちょうする (seichō suru)}]{\ja{せいちょうする} (seichō suru) — to grow up; to grow}

\label{lesson:JA-C287-seichou-suru}



\begin{quote}
Before the new one: say the Japanese for a grandmother, then the Japanese for a relative.
\end{quote}

\subsection*{You'll want to know: \ja{せいちょうする}}
\textbf{\ja{せいちょうする}} — \emph{seichō suru} — \textquotedblleft{}to grow up; to grow\textquotedblright{}.

\begin{quote}
\textbf{\ja{こどもは} \ja{すぐ} \ja{せいちょうします。}} — \emph{kodomo wa sugu seichō shimasu} — \textquotedblleft{}Children grow up fast.\textquotedblright{}
\end{quote}

Polite \textbf{\ja{せいちょうします}}, \emph{seichōshimasu}; a \ja{する} verb. Past (polite): \textbf{\ja{こどもが} \ja{せいちょうしました。}}, \emph{kodomo ga seichō shimashita}.

\begin{cousinweb}
\textbf{\ja{せいちょうする}} (\emph{seichō suru}) is Sino-Japanese (kango): two Chinese characters read with their on'yomi, readings borrowed from Chinese. They are \emph{sei}, \textquotedblleft{}grow\textquotedblright{} (Middle Chinese \emph{dzyeng}) and \emph{chō}, \textquotedblleft{}long\textquotedblright{} (\emph{trjangX}) (KANJIDIC2; Middle Chinese in Baxter's transcription). The native verb \emph{suru}, \textquotedblleft{}to do\textquotedblright{}, makes it a verb.
\end{cousinweb}

\begin{grammarlens}[title={What you've built}]
One more word for telling what happened.
\end{grammarlens}

\subsection*{Your turn}
\begin{itemize}
\raggedright
  \item \emph{Say it:} \emph{seichō suru}
  \item \emph{Say it:} \emph{seichō suru}, once more
  \item \emph{Say it:} \emph{kodomo wa sugu seichō shimasu} — Children grow up fast.
  \item \emph{Recall:} say the Japanese for a grandmother, then the Japanese for a relative, then say \emph{seichō suru} again
\end{itemize}

\begin{tcolorbox}[breakable,colback=teal!4,colframe=teal!35!black,title={Before you move on}]
What does \ja{せいちょうする} mean? (\textquotedblleft{}To grow up; to grow\textquotedblright{}.) Say it once more.
\end{tcolorbox}
